% =====================================================================
%  CHAPTER 1 - INTRODUCTION
% =====================================================================
\clearpage
\chapter{Introduction}

\section{Background}

Much of the software that runs a business is not a single request answered in
milliseconds but a process that unfolds over seconds, hours or days. Placing a
food order charges a card, reserves stock, notifies a restaurant and assigns a
driver. Settling an invoice debits one ledger, credits another and files a tax
record. An agent built on a large language model that is asked to arrange
travel plans an itinerary, calls a booking service, waits for a person to
approve the result and then confirms it. Each of these is a \emph{workflow}: a
sequence of steps, spread over several independent services, in which every
step depends on the results of the steps before it and most steps change the
state of some system outside the workflow itself.

Two properties make workflows hard to run reliably. A workflow is stateful,
because step four only makes sense in the light of what steps one to three
decided. Its steps also have external side effects that the workflow cannot
withdraw: once a payment provider has moved money, no local rollback can move
it back. A database transaction has neither difficulty, since it is short
enough to hold its locks throughout and its effects are confined to a system
that can undo them. A workflow has both.

As a result, an ordinary crash becomes a correctness problem. Suppose the
order workflow is interrupted after the card has been charged and the stock
reserved, but before the restaurant has been told. A system that has kept no
durable record of its progress has two choices when it restarts. It can run
the workflow again from the beginning, charging the card a second time, or it
can abandon the workflow, keeping the customer's money for an order that will
never be delivered. Without a record of what was already done, there is no
correct third option.

The usual responses in practice are retry loops, periodic jobs that poll a
status column, and message queues whose consumers record progress in
application tables. Each works for the workflow it was written for, and each
must be written again for the next one. \emph{Durable execution} is the name
given to systems that solve the problem once for every workflow built on
them~\cite{temporaldocs, restatedocs}. The idea on which they rest, and on
which this report rests, is that \emph{every decision a workflow makes is
appended to a log before it is acted on, and recovery is performed by running
the workflow's own code again against that log.} Completed steps return their
recorded results instead of executing a second time, so the re-executed code
arrives at the point the failed process had reached. Because nothing private
to one machine is part of the workflow's state, the recovery can take place on
a different machine.

Table~\ref{tab:terminology} fixes the terms used throughout the report.

\begin{table}[htbp]
  \centering
  \caption{Terminology used in this report}
  \label{tab:terminology}
  \small
  \begin{tabular}{@{}p{1.15in}p{4.4in}@{}}
    \toprule
    \textbf{Term} & \textbf{Meaning} \\
    \midrule
    Workflow & Orchestration code that decides what happens next and performs
      no input or output itself. It must be deterministic, because it is
      re-executed on every recovery. \\
    Activity & One step with an external side effect, such as charging a card
      or calling an API. It may be non-deterministic and runs at least once. \\
    Event & An immutable record appended to the log describing something that
      has already happened. \\
    History & The ordered events of one workflow run; the complete state of
      that run. \\
    Replay & Re-executing a workflow against its history so that completed
      steps resolve from the log. \\
    Shard & One of a fixed number of partitions. A workflow belongs to one
      shard for its whole life. \\
    Lease & A time-limited claim on a shard, held in a database row and
      renewed periodically, that makes one node the owner of the shard. \\
    Fencing token & A counter that increases whenever ownership of a shard
      changes; a write is accepted only with the current value. \\
    Saga & A multi-step transaction that, if it fails part way, is undone by
      running a compensating action for each step that had succeeded. \\
    \bottomrule
  \end{tabular}
\end{table}

Industry has adopted this approach widely. Uber's Cadence and its successor
Temporal run very large numbers of workflows for ride and delivery
systems~\cite{temporaldocs}; Microsoft's Durable Functions provide the same
model for serverless code~\cite{durablefunctions2021}; and the DBOS project
argues that durable execution belongs inside the database rather than in a
separate service~\cite{dbos2022, r3vldb2023}. The internal mechanisms of such
an engine, namely how the log is partitioned, how state is reconstructed
deterministically and how ownership of a partition moves safely between
machines when one fails, are less often built from first principles. This
project builds one, and evaluates it on a cloud cluster under load and under
injected failures.

\section{Motivation}

\subsection{Existing engines depend on several services}

The production engines surveyed for this work each require four to six
cooperating infrastructure components. Netflix's Conductor uses a database,
a search cluster, a queue, a server tier and a worker fleet. Temporal runs
four internal services over a database and, for visibility queries, a search
cluster~\cite{temporaldocs}. Meta's Shard Manager, which assigns partitions
across very large fleets, relies on ZooKeeper for coordination~\cite{shardmanager2021,
zookeeper2010}. Each additional component is a separate failure domain that
must be configured, monitored and understood during an incident, and for a
small deployment the orchestration system can cost more to operate than the
application it serves. This project asks whether the same guarantees can be
provided with a single dependency.

\subsection{Partitioning and simplicity are usually traded off}

The simple systems are not partitioned and the partitioned systems are not
simple. DBOS obtains durable execution with very few moving parts but runs on
one node~\cite{dbos2022}; Temporal and Shard Manager partition work across
machines but add a coordination service or a multi-service control plane to
decide which machine owns what. Our position is that the two need not be in
tension, because a relational database already provides the linearizable
compare-and-swap that partition ownership requires~\cite{leases1989,
kleppmann2017}. What is needed is a careful use of one row per shard rather
than a consensus cluster~\cite{raft2014}.

\subsection{Agents defeat byte-level deduplication}

Durable execution engines prevent duplicate side effects by recognising a
repeated request, and they recognise it by comparing bytes: a hash, an
idempotency key or an exact match on recorded arguments. This is sound when
the arguments are fixed at the moment a step is first scheduled. It is not
sound when the caller is a language-model agent, which after a crash
\emph{re-plans} rather than replays: an instruction recorded as
``BLR$\rightarrow$DEL Friday morning'' may be issued again as ``Bangalore to
Delhi, Friday AM''. The two strings hash differently and the flight is booked
twice. Step repetition of this kind appears as a recurring category in
empirical studies of multi-agent failures~\cite{whyfail2025} and in
Microsoft's taxonomy of agentic failure modes~\cite{msredteam2025}.

\section{Problem Statement}

The problem addressed in this work is to design and implement a horizontally
sharded durable execution engine that runs long, multi-step workflows to
completion despite process crashes, node failures and partial failures, using
event sourcing, deterministic replay, idempotent activity execution and
lease-based ownership; that is deployable on a small cluster with PostgreSQL
as its only infrastructure dependency; that detects semantically equivalent
repeated actions issued by a recovering agent before the side effect occurs a
second time; and to measure its throughput, its sensitivity to its main
configuration parameters, and its recovery from injected faults on real cloud
infrastructure.

Two limits are part of the statement. The engine targets correctness and fault
tolerance at a moderate, measurable scale rather than the throughput of
large production deployments. Ownership is obtained from a lease in a
linearizable database rather than from a consensus protocol, which makes the
database the availability floor of the system.

\section{Objectives}

The work has eight objectives, listed in Table~\ref{tab:objectives} with the
mechanism that realises each and the chapter in which it is evaluated.

\begin{table}[htbp]
  \centering
  \caption{Project objectives}
  \label{tab:objectives}
  \small
  \begin{tabular}{@{}p{0.2in}p{2.75in}p{1.65in}p{0.75in}@{}}
    \toprule
    \textbf{No.} & \textbf{Objective} & \textbf{Mechanism} & \textbf{Evaluated} \\
    \midrule
    1 & Record every state transition as an immutable event so that a run can
        be reconstructed exactly & Append-only event log & Ch.~6, 7 \\
    2 & Resume a failed run from its history rather than from the start &
        Deterministic replay & Ch.~6, 7 \\
    3 & Never repeat an external side effect across retries & Stable
        deduplication keys & Ch.~6 \\
    4 & Partition execution across nodes & Fixed shards, consistent-hash
        ring & Ch.~7 \\
    5 & Lose no progress when a shard owner fails & Leases with fencing
        tokens & Ch.~7 \\
    6 & Undo multi-step transactions that fail part way & Saga
        compensation & Ch.~6 \\
    7 & Detect reworded duplicate agent actions & Journalled semantic
        verdict & Ch.~6 \\
    8 & Measure capacity, sensitivity and recovery on cloud hardware & EKS
        test bed, fault injector & Ch.~7 \\
    \bottomrule
  \end{tabular}
\end{table}

\section{Scope and Limitations}

The scope is the engine, two demonstration applications built on it, and the
test bed used to evaluate it. The main limitations are stated here so that the
design and the results can be read with them in mind.

\begin{enumerate}[leftmargin=*, itemsep=2pt]
  \item The engine provides at-least-once execution of activities with strong
    deduplication, not exactly-once execution. A crash after a side effect has
    happened but before it has been recorded cannot be told apart from a crash
    before the side effect. The engine supplies a stable deduplication key
    that the downstream service can use to close this window.
  \item Ownership relies on PostgreSQL as the single source of truth. The
    cloud evaluation used one database host with no standby, so it says
    nothing about database high availability.
  \item The event log is never garbage collected; retention is future work.
  \item Semantic deduplication is opt-in and conservative. A false negative is
    no worse than exact matching, but a false positive would suppress an
    action the user wanted, so the semantic mode must be requested per
    activity and every verdict is journalled.
  \item All cloud measurements come from one availability zone on burstable
    (``flex'') instance types, over individual runs of a few minutes. They
    establish the shape of the behaviour on this configuration, not absolute
    capacity for other hardware.
\end{enumerate}

\section{Contributions}

The contributions of this work are as follows.

\begin{enumerate}[leftmargin=*, itemsep=2pt]
  \item A durable execution engine in which event sourcing, replay, a durable
    task queue, timers, signals, sagas and shard failover are implemented over
    PostgreSQL alone, with no message broker or coordination service.
  \item Fencing at the persistence layer: every transaction that writes shard
    state re-reads the shard's fencing token with \cd{SELECT ... FOR UPDATE}
    and commits only if the token is still the writer's. On the cloud test bed
    this rejected every write attempted by a stale owner after a process
    freeze.
  \item A two-level mapping in which workflows hash to a fixed set of shards
    and shards are assigned to nodes by a consistent-hash ring, so that a
    change of membership moves ownership but no data.
  \item A semantic idempotency check that compares the meaning of an agent's
    intended action with the actions already in the run's history and
    journals the verdict, keeping replay deterministic and auditable.
  \item A cloud evaluation on Amazon EKS covering idle overhead, latency
    floor, closed- and open-loop capacity, sensitivity to workers, shards and
    replicas, and recovery from six classes of injected fault, together with
    the evaluation harness and its validity checks.
\end{enumerate}

\section{Organisation of the Report}

Chapter~2 reviews related work on durable execution, partition management,
ownership and fencing, sagas, and semantic similarity, and identifies the gap
this work addresses. Chapter~3 states the requirements and the analysis that
led to the chosen approach. Chapter~4 presents the high-level and low-level
design. Chapter~5 describes the implementation and Chapter~6 the testing
methodology, including the cloud test bed and the fault-injection procedure.
Chapter~7 reports and discusses the results, and Chapter~8 concludes and
outlines future work.
